\documentclass[a4paper, 12pt, notitlepage]{report}

\usepackage{amsfonts}  
\usepackage{graphicx}  
\usepackage{fullpage}
\usepackage[pagewise]{lineno}
\usepackage{hyperref}
\usepackage{longtable}                      
\usepackage{ltxtable}                       
\usepackage[usenames,dvipsnames]{color}
\usepackage{colortbl}          
\usepackage{tabularx, booktabs}   
\usepackage{multirow}
\usepackage{dirtree}
\definecolor{tableheadcolor}{gray}{0.8}
\definecolor{darkblue}{rgb}{0,0,.5}
\definecolor{gray}{gray}{0.95}
\definecolor{lightgray}{gray}{0.95}
\newcommand{\file}[1]{\textbf{#1}}
\newcommand{\state}[1]{{\tt\textbf{#1}}}
\newcommand{\signal}[1]{{\textbf{\tt#1}}}
\title{YAC - Yet Another CORDIC Core} 
\author{Christian Haettich [feddischson@gmail.com]} 
\date{\today}

\begin{document}

%%%%%%%%%% PRELIMINARY MATERIAL %%%%%%%%%%
\maketitle
\begin{center}
IP core and C-model documentation 
\end{center}
\thispagestyle{empty}

\tableofcontents 


\chapter{Organization and Introduction}
\section{History}
\begin{table}[htbp]
   \center
   \begin{tabular}{@{}llll@{}}
      \rowcolor{tableheadcolor}\textbf{Version}& \textbf{Date} & \textbf{Description} & \textbf{Author}  \\

      \multirow{1}{*}{0.02}  & 26th June 2018 & E-mail address update and diagrams de-drawn & Christian Haettich\\\midrule
      \multirow{1}{*}{0.01}  & 4th March 2014 & Initial document & Christian Haettich\\\midrule
      \bottomrule
   \end{tabular}
   \caption{History}
   \label{tab:frqband}
\end{table} 


\section{Folder and Files}

\dirtree{%
.1 /.
.2 c\_model.
.3 cordic\_pipeline.c\DTcomment{C-code of the bit accurate model of the pipelined core}.
.3 cordic\_pipeline.h\DTcomment{C header of the pipelined model}.
.3 cordic\_pipeline.py\DTcomment{Python (ctypes) wrapper of the pipelined model}.
.3 cordic\_pipeline\_tables.h\DTcomment{Generated tables, do not edit}.
.3 gen\_tables.py\DTcomment{Generates cordic\_pipeline\_tables.h and cordic\_pipeline\_pkg.vhd}.
.2 c\_octave.
.3 cordic\_iterative.c\DTcomment{C-code for bit accurate model}.
.3 cordic\_iterative\_code.m\DTcomment{Script to auto-generate VHDL and C-code}.
.3 cordic\_iterative\_test.m\DTcomment{YAC performance analysis script}.
.2 doc.\DTcomment{Contains files to create this documentation}.
.2 licenses.
.3 lgpl-3.0.txt\DTcomment{LGPL license}.
.2 README.md\DTcomment{A short read-me file}.
.2 rtl.
.3 vhdl.
.4 cordic\_pipeline\_int.vhd\DTcomment{Pipelined core}.
.4 cordic\_pipeline\_pkg.vhd\DTcomment{Pipelined core package, generated}.
.4 tb\_cordic\_pipeline.vhd\DTcomment{VUnit testbench of the pipelined core}.
.4 cordic\_iterative\_int.vhd\DTcomment{Top-level VHDL file}.
.4 cordic\_iterative\_pkg.vhd\DTcomment{VHDL package file}.
.4 cordic\_iterative\_tb.vhd\DTcomment{VHDL testbench}.
.2 test.
.3 cordic\_test\_common.py\DTcomment{Shared helpers of the pipelined tests}.
.3 model.
.4 run.py\DTcomment{C-model against the exact rotation and vectoring, writes the sweeps}.
.3 pipeline.
.4 run.py\DTcomment{VUnit, RTL against the C-model}.
}

\section{License}
Copyright (c) 2026, Christian Haettich, All rights reserved. \newline\newline 
YAC is free software; you can redistribute it and/or               
modify it under the terms of the GNU Lesser General Public         
License as published by the Free Software Foundation; either       
version 3.0 of the License, or (at your option) any later version. \newline\newline
YAC is distributed in the hope that it will be useful,             
but WITHOUT ANY WARRANTY; without even the implied warranty of     
MERCHANTABILITY or FITNESS FOR A PARTICULAR PURPOSE.  See the GNU  
Lesser General Public License for more details.              \newline\newline
You should have received a copy of the GNU Lesser General Public   
License along with this library. If not, download it from          
http://www.gnu.org/licenses/lgpl                                   


\section{The CORDIC algorithm}

The CORDIC algorithm is used to do trigonometric, hyperbolic, multiplication 
and division calculations in a hardware-efficient way. This means, that
only bit-shift, addition and subtraction operations in combination
with a lookup-table is required.\newline
Good introductions are available in \cite{survey}\cite{vlsi}\cite{dawid}.
A description on wikibook.org \cite{wikibook} is also useful.
For this reason, no more introduction is given here and the following chapters assume, that
the reader is familiar with the CORDIC algorithm.


\chapter{IP-Core Description}

The two files \file{cordic\_iterative\_int.vhd} and \file{cordic\_iterative\_pkg.vhd}
implements an iterative CORDIC algorithm in fixed-point format. 
Iterative means, that the IP-core is started with a set of input data and after
a specific amount of clock cycles, the result is available. No parallel data
can be processed by this core; the pipelined core of chapter \ref{chap:pipeline}
does that. The 
following six modes are supported:
\begin{itemize}
   \item trigonometric rotation
   \item trigonometric vectoring
   \item linear rotation
   \item linear vectoring
   \item hyperbolic rotation
   \item hyperbolic vectoring
\end{itemize}

\section{Port Description}
\begin{table}[htbp]
   \center
   \begin{tabular}{@{}lllll@{}}
      \rowcolor{tableheadcolor}\textbf{Name}& \textbf{Type} & \textbf{Direction} & \textbf{Size} & \textbf{Description} \\

      \multirow{1}{*}{clk}    & std\_logic         & input  & 1 & clock              \\\midrule
      \multirow{1}{*}{rst}    & std\_logic         & input  & 1 & synchronous reset  \\\midrule
      \multirow{1}{*}{en}     & std\_logic         & input  & 1 & clock-enable       \\\midrule
      \multirow{1}{*}{start}  & std\_logic         & input  & 1 & start trigger      \\\midrule
      \multirow{1}{*}{done}   & std\_logic         & output & 1 & done indicator     \\\midrule
      \multirow{1}{*}{mode}   & std\_logic\_vector & input  & 4 & mode configuration \\\midrule
      \multirow{1}{*}{x\_i}   & std\_logic\_vector & input  & XY\_WIDTH                & X input vector        \\\midrule
      \multirow{1}{*}{y\_i}   & std\_logic\_vector & input  & XY\_WIDTH                & Y input vector        \\\midrule
      \multirow{1}{*}{a\_i}   & std\_logic\_vector & input  & A\_WIDTH + 2             & angular input vector  \\\midrule
      \multirow{1}{*}{x\_o}   & std\_logic\_vector & input  & XY\_WIDTH + GUARD\_BITS  & X output vector       \\\midrule
      \multirow{1}{*}{y\_o}   & std\_logic\_vector & input  & XY\_WIDTH + GUARD\_BITS  & Y output vector       \\\midrule
      \multirow{1}{*}{a\_o}   & std\_logic\_vector & input  & A\_WIDTH + 2             & angular output vector \\\midrule
      \bottomrule
   \end{tabular}
   \caption{Port description}
   \label{tab:ports}
\end{table} 

\begin{figure}[p]
\centering
   \includegraphics[width=250pt]{figs/block_symbol.pdf}
   \caption{YAC block symbol}
   \label{fig:ports}
\end{figure}
Table \ref{tab:ports} and Figure \ref{fig:ports} gives an overview
of the YAC ports.
The three input ports start, done and mode and are used to interact
with an external state machine or with a software driver.
After setting \signal{start} to high, the YAC starts processing.
If all necessary rotations are done, the \signal{done} output is set to 
high for one clock cycle.
When starting the YAC, all other inputs must contain valid data.
The mode input is used to select the CORDIC mode.
Table \ref{fig:ports} shows all input and output ports and Figure 
\ref{fig:ports} shows a CORDIC block symbol: $x_i, y_i, a_i, x_o, y_o, a_o$
are data ports whereas start, done and mode are configuration signals.






\section{Parameter Description}
\begin{table}[htbp]
   \center
   \begin{tabular}{@{}lllll@{}}
      \rowcolor{tableheadcolor}\textbf{Name}& \textbf{type} & \textbf{size} & \textbf{Description} \\

      \multirow{1}{*}{XY\_WIDTH}    & natural  & defines the size of x and y input and output vectors \\\midrule
      \multirow{1}{*}{A\_WIDTH}     & natural  & defines the size of the angular input and output vectors \\\midrule
      \multirow{1}{*}{GUARD\_BITS}  & natural  & defines the number of guard bits \\\midrule
      \multirow{1}{*}{RM\_GAIN}     & natural  & defines the precision of the CORDIC gain removal \\\midrule
      \bottomrule
   \end{tabular}
   \caption{Parameter description}
   \label{tab:params}
\end{table} 
Table \ref{tab:params} shows the four parameter, which can be used to parameterize 
the YAC.


\section{Mode description}
With the mode-input, it is possible to select between circular, linear and hyperbolic 
operation as well as between rotation and vectoring operations:
\begin{itemize}
\item mode(1:0) = "00" selects circular operation
\item mode(1:0) = "01" selects linear operation
\item mode(1:0) = "10" selects hyperbolic operation
\item mode(1:0) = "11" is reserved
\item mode(3) = '0' selects rotation
\item mode(3) = '1' selects vectoring
\end{itemize}
The package cordic\_iterative\_int.vhd contains the following defines to
setup the mode:
\begin{itemize}
   \item VAL\_MODE\_CIRC = "00" 
   \item VAL\_MODE\_LIN  = "01" 
   \item VAL\_MODE\_HYP  = "10" 
   \item I\_FLAG\_VEC\_ROT = 3 (bit index)
\end{itemize}
For example if a hyperbolic rotation mode is required, the mode input is "0010", 
and for the linear vector operation, the mode input is "1001".
Please note, that bit number 2 is reserved for future implementations.\newline
Table \ref{tab:modes} defines the behaviour of the YAC 
for different mode configurations.
\begin{table}[htbp]
   \center
   \begin{tabular}{@{}ll@{}}
      \rowcolor{tableheadcolor}\textbf{Setup}& \textbf{Description} \\

      \multirow{2}{*}{mode =( FLAG\_VEC\_ROT = 1, VAL\_MODE\_CIR ) } & $a_o = \textrm{atan2}( y_i, x_i ) \cdot \beta,  $\\
                                                                     & $x_o = \sqrt{ x_i^2+y_i^2 }  $ \\\midrule 
      \multirow{2}{*}{mode =( FLAG\_VEC\_ROT = 0, VAL\_MODE\_CIR ) } & $x_o = \cos( a_i / \beta ) \cdot \alpha  $\\
                                                                     & $y_o = \sin( a_i / \beta ) \cdot \alpha  $\\\midrule 
      \multirow{1}{*}{mode =( FLAG\_VEC\_ROT = 1, VAL\_MODE\_LIN ) } & $a_o = z + y/x $\\\midrule
      \multirow{1}{*}{mode =( FLAG\_VEC\_ROT = 0, VAL\_MODE\_LIN ) } & $a_o = y + x\cdot z  $ $  $\\\midrule
      \multirow{2}{*}{mode =( FLAG\_VEC\_ROT = 1, VAL\_MODE\_HYP ) } & $a_o = \textrm{tanh}( y_i / x_i ) \cdot \beta,   $\\
                                                                     & $x_o = \sqrt{ x_i^2-y_i^2 }                                 $\\\midrule 
      \multirow{2}{*}{mode =( FLAG\_VEC\_ROT = 0, VAL\_MODE\_HYP ) } & $x_o = \cosh( a_i / \beta ) \cdot \alpha$\\
                                                                     & $y_o = \sinh( a_i / \beta ) \cdot \alpha$\\
      \bottomrule
   \end{tabular}
   \caption{Mode description (see equation \ref{eqn:alpha} and \ref{eqn:beta} for $\alpha$ and $\beta$)}
   \label{tab:modes}
\end{table} 


\section{Data Range and Limitations}

   Both, the $x_i$ and $y_i$ inputs are defined with a size of $XY\_WIDTH$ and therefore,
   the maximum positive value in the two's complement format is
   \begin{eqnarray}
   \alpha & = & 2^{XY\_WIDTH-1}-1
   \label{eqn:alpha}
   \end{eqnarray}
   The angle input $a_i$ is defined with a size of $A\_WIDTH+2$.
   The value $\beta$ is defined with
   \begin{eqnarray}
   \beta & = & 2^{A\_WIDTH-1} -1
   \label{eqn:beta}
   \end{eqnarray}

   \subsection{Circular Vectoring Mode}

      The valid input range for $x_i$ and $y_i$ is $-\alpha ... +\alpha$.
      The angular input $a_i$ is ignored.

   \subsection{Circular Rotation Mode}
      The valid input range for $x_i$ and $y_i$ is $-\alpha ... +\alpha$.
      For the angular input $a_i$, values between $-\beta\pi$ and $+\beta\pi$ are valid.
      For calculating a complex rotation of a complex vector, all three inputs
      are required. For calculating $\sin$ and $\cos$, $x_i$ is set to $\alpha$
      and $y_i$ to 0, the angle input $a_i$ gets the angle.

   \subsection{Linear Vectoring Mode}
      The linear vectoring mode is used to calculate $y_i / x_i$. The limitation
      of this operation is the following:
      \begin{eqnarray}
         \frac{y_i}{x_i} \le 2
      \end{eqnarray}
      The valid input values for $x_i$ and $y_i$ are $-\alpha ... +\alpha$.

   \subsection{Linear Rotation Mode}
      The two inputs $x_i$ and $y_i$ have a valid data range between $-\alpha$ and $+\alpha$.


   \subsection{Hyperbolic Vectoring Mode}
      The data-range for $x_i$ and $y_i$ is $-0.79 \cdot \alpha$ to $0.79 \cdot \alpha$.


   \subsection{Hyperbolic Rotation Mode}
      The data-range for $a_i$ is $-\beta ... \beta$. Typically, $x_i$ is set to 
      $\alpha$ and $y_i$ to 0.

\section{Internal Operation}
\begin{figure}[p]
\centering
   \includegraphics[width=300pt]{figs/state_diagram.pdf}
   \caption{State diagram}
   \label{fig:states}
\end{figure}
Five states are used do to the CORDIC algorithm. The five states are visualized in a
state diagram in Figure \ref{fig:states}.
After reset, the YAC goes into the \state{ST\_IDLE} state. Only in this state, the 
start signal is accepted. After starting the YAC, the state switches from
the \state{ST\_IDLE} state to the \state{ST\_INIT} state. The init state does some initial 
rotations/flipping.
After initialization, the main rotation starts with the \state{ST\_ROTATION} state
(there are a few cases, where no further rotations are required, see next sections).
Every rotation step is done within two clock cycles: in the first cycle, 
the angular value is loaded from the lookup-table and shifted. In addition, 
the x and y values are shifted. In the second
clock cycle, the results are added according to the CORDIC algorithm.
After all required iterations are done, the state switches to
the \state{RM\_GAIN} state, where the CORDIC gain is removed (depending on 
the mode). The last state \state{ST\_DONE} is only used to set the done flag to 
'1' for one clock cycle. After this, the YAC returns to the \state{ST\_IDLE} state.



   \subsection{Initialization}
   During the initialization state ST\_INIT, pre-rotations are done, depending on 
   the selected mode.


      \subsubsection{Circular Vectoring Mode}
      Because we support atan2 and not atan, we do the following pre-rotations. 
      Some specific input values don't need a further processing, for 
      example $atan( 0, 0 )$.\newline\newline
         \begin{tabular}{@{}lp{380pt}@{}}
            \rowcolor{tableheadcolor}\textbf{Input} &\textbf{ Description }  \\

            $x_i = 0, y_i = 0$    &    Fixed result: $x_o =  0, a_o = 0$,   
            to support cartesian to polar coordinate transformations, 
            further processing is skipped  \\ \midrule

            $x_i = 0, y_i > 0 $   &    Fixed result: $x_o =  y_i, a_o = \pi/2$,  
            further processing is skipped \\ \midrule

            $x_i = 0, y_i < 0 $   &    Fixed result: $x_o = -y_i, a_o = -\pi/2$,  
            further processing is skipped \\ \midrule


            $x_i < 0, y_i < 0 $   &    Pre-rotation from the third to the first quadrant,
                                       the angular register is initialized with $-\pi$ and
                                       the sign of the $x$ and $y$ register is flipped. \\ \midrule

            $x_i < 0, y_i \ge 0 $   &    Pre-rotation from the second to the fourth quadrant,
                                       the angular register is initialized with $\pi$ and
                                       the sign of the $x$ and $y$ register is flipped. \\ \midrule


            $x_i = + \alpha, y_i = + \alpha $   &  Fixed result: $x_o = \sqrt{2} \cdot \alpha, a_o = \frac{\pi}{4} \cdot \beta$ ,
            further processing is skipped \\ \midrule

            $x_i = + \alpha, y_i = - \alpha $   &  Fixed result: $x_o = \sqrt{2} \cdot \alpha, a_o = - \frac{\pi}{4} \cdot \beta$, 
            further processing is skipped \\ \midrule

            $x_i = - \alpha, y_i = + \alpha $   &  Fixed result: $x_o = \sqrt{2} \cdot \alpha, a_o =  \frac{3 \pi}{4} \cdot \beta$, 
            further processing is skipped \\ \midrule

            $x_i = - \alpha, y_i = - \alpha $   &  Fixed result: $x_o = \sqrt{2} \cdot \alpha, a_o = - \frac{3 \pi }{4} \cdot \beta$ ,
            further processing is skipped \\ 
            \bottomrule
         \end{tabular}

      \subsubsection{Circular Rotation Mode}
      The following pre-rotations are done:\newline\newline
         \begin{tabular}{@{}lp{380pt}@{}}
            \rowcolor{tableheadcolor}\textbf{Input} &\textbf{ Description }  \\

            $a_i < - \frac{\pi}{2}$    & Pre-rotation from the third to the first quadrant, 
                                         the sign of the $x$ and $y$ register is flipped, 
                                         the angular register is initialized with $\pi$   \\ \midrule

            $a_i >   \frac{\pi}{2}$    & Pre-rotation from the second to the fourth quadrant, 
                                         the sign of the $x$ and $y$ register is flipped, 
                                         the angular register is initialized with $-\pi$   \\
            \bottomrule
         \end{tabular}

      \subsubsection{Linear Vectoring Mode}
      The following pre-rotations are done:\newline\newline
         \begin{tabular}{@{}lp{380pt}@{}}
            \rowcolor{tableheadcolor}\textbf{Input} &\textbf{ Description }  \\

            $x_i < 0$          & Pre-rotation from the left half to the right half, thus
                                 the sign of the $x$ and $y$ register is flipped \\
            \bottomrule
         \end{tabular}


   \subsection{Rotation}

   \subsubsection{Hyperbolic Repetitions}
   The following iteration steps are repeated for convergence of the hyperbolic mode:
   \begin{eqnarray}
      4, 13, 40, 121, ..., 3i+1
   \end{eqnarray}



\section{Testbench}
A VHDL testbench (cordic\_iterative\_tb.vhd) is available which reads from an input file
data. This data is fed into the YAC. In addition, the input file must
provide the output data, which is expected from the YAC processing.
The following format is used in the testbench input-file:
\begin{verbatim}
x_i     y_i    a_i   x_o   y_o   a_o   mode

1234   1234   1234  1234  1234  1234    111  \
2345   2345   2345  2345  2345  2345    222  |  7 columns:
3456   3456   3456  3456  3456  3456    333  |  3 input, 3 output and one mode
...    ...                                   |
...                                          |
...
\end{verbatim}
The data must be in the two's complement with a base of 10.
After processing all lines, the testbench prints some info, how much
tests failed.\newline
For generating the test-patterns, the C-model with octave is used. By running
the script \file{cordic\_iterative\_test.m}, an entry 
in the testbech data file is done for every cordic calculation.


\chapter{Using the C-Model with Octave or Matlab}
For using the C-model, open octave or Matlab and compile the C-Model
with \newline
\textit{mex cordic\_iterative.c}
Afterwards, the C-model can be used in Octave/Matlab with the function call\newline
{\tt [ x\_o, y\_o, a\_o, it ] = cordic\_iterative( x\_i, y\_i, a\_i, mode, XY\_WIDTH, ANGLE\_WIDTH,
GUARD\_BITS, RM\_GAIN )}.\newline
The input and output arguments reflect almost the ports and parameters of the 
VHDL toplevel implementation, except that there is no done, start clk and rst port 
in the C-model. In addition, the last argument provides the number of iterations,
which are done.\newline \newline
The input arguments x\_i, y\_i, and a\_i can be 1xN vectors with N elements.
In this case, this three input arguments must have the same size. All
output arguments will also have the same size 1xN.



\section{Guard-Bits -- Input and output sizes}
It is possible to define MSB guard bits.
The need of MSB guard bits have several reasons. On the one hand,
they are required to handle the data growth of the cordic gain within the rotation state.
(which is compensated later). A nother reason is the data growth caused by the choise of
input data. For example 
\begin{eqnarray}
   \sqrt{ \alpha^2 + \alpha^2 } = \sqrt{2} \alpha
\end{eqnarray}
and therefore, a growth by the factor $\sqrt{2}$ happens, if an absolute calculation
of the two maximum input values is processed.\newline
\newline
Another point is the input and output size of the angular values ($a_i$ and $a_o$)
Because within this YAC, $\beta$ represents the angle 1, the input and output with
of $a_i$ and $a_o$ is A\_WIDTH+2, which allows to process angle values within
$-\pi ... \pi$


\subsection{Input data width versus iterations}
The number of iterations depends on the width of the input data. 
Input data with $X$ bit require $X+1$ rotations \cite{dawid}.
This requires, that the number of iterations are adapted. One solution might
be to detect the size of the input data and use this information for the CORDIC processing.
In YAC, a different sheme was chosen: the rotations are done until 
\begin{itemize}
   \item the angular register is 0 or doesn't change (in case of the rotation mode) 
   \item the y register is 0 or doesn't change (in case of the vecotring mode).
\end{itemize}
This results in a dynamic iterations adaption, depending the the input data width.

\chapter{Pipelined Core}
\label{chap:pipeline}
\sloppy

The entity \file{cordic\_pipeline\_int} in \file{cordic\_pipeline\_int.vhd} is an unrolled
CORDIC for circular rotation or, selected by the generic \signal{VECTORING}, circular vectoring. It accepts a new sample every clock cycle and has a fixed
latency. It uses \file{numeric\_std} only, with no vendor primitives. The package
\file{cordic\_pipeline\_pkg.vhd} holds the angle table, the gain-removal terms and the
function \signal{cordic\_pipeline\_latency( N\_ITER, RM\_GAIN )}.
The iterative core and its package are not used by the pipelined core.

\section{Port Description}
\begin{table}[htbp]
   \center\small
   \begin{tabular}{@{}llllp{90pt}@{}}
      \rowcolor{tableheadcolor}\textbf{Name}& \textbf{Type} & \textbf{Direction} & \textbf{Size} & \textbf{Description} \\

      clk     & std\_logic         & input  & 1 & clock, rising edge \\\midrule
      rst     & std\_logic         & input  & 1 & synchronous reset, active high \\\midrule
      en      & std\_logic         & input  & 1 & clock-enable, freezes the whole pipeline \\\midrule
      valid\_i & std\_logic        & input  & 1 & $x_i, y_i, a_i$ carry a sample \\\midrule
      x\_i    & std\_logic\_vector & input  & XY\_WIDTH                & real part, two's complement \\\midrule
      y\_i    & std\_logic\_vector & input  & XY\_WIDTH                & imaginary part, two's complement \\\midrule
      a\_i    & std\_logic\_vector & input  & A\_WIDTH                 & binary angle, ignored in vectoring \\\midrule
      valid\_o & std\_logic        & output & 1 & $x_o, y_o, a_o$ carry a result \\\midrule
      x\_o    & std\_logic\_vector & output & XY\_WIDTH + GUARD\_BITS  & real part of the result \\\midrule
      y\_o    & std\_logic\_vector & output & XY\_WIDTH + GUARD\_BITS  & imaginary part of the result \\\midrule
      a\_o    & std\_logic\_vector & output & A\_WIDTH                 & rotation: all zeros; vectoring: angle of $x_i + i\,y_i$ \\
      \bottomrule
   \end{tabular}
   \caption{Port description of the pipelined core}
   \label{tab:pipe_ports}
\end{table}

\section{Parameter Description}
The entity asserts the ranges with severity failure.
\begin{table}[htbp]
   \center\small
   \begin{tabular}{@{}lllp{200pt}@{}}
      \rowcolor{tableheadcolor}\textbf{Name}& \textbf{Default} & \textbf{Range} & \textbf{Description} \\

      XY\_WIDTH   & 16 & $\ge 2$  & width of $x_i$ and $y_i$ \\\midrule
      A\_WIDTH    & 18 & 3 to 32  & angle width, $2^{A\_WIDTH}$ is a full circle \\\midrule
      GUARD\_BITS & 2  & $\ge 0$  & MSB guard bits \\\midrule
      LSB\_BITS   & 0  & $\ge 0$  & internal fraction bits, dropped at the output by truncation \\\midrule
      N\_ITER     & 16 & 1 to 32  & CORDIC iterations \\\midrule
      RM\_GAIN    & 4  & 0 to 12  & gain-removal terms, 0 means no removal \\\midrule
      VECTORING   & false & boolean & false: rotation, true: vectoring \\
      \bottomrule
   \end{tabular}
   \caption{Parameters of the pipelined core}
   \label{tab:pipe_params}
\end{table}
In addition, $XY\_WIDTH + GUARD\_BITS + LSB\_BITS$ is at most 62. The angle table has
32 entries of 32 bits ($\mathrm{TABLE\_WIDTH} = 32$); the core uses the top $A\_WIDTH$ bits of
each entry. The default values are placeholders: the FFT accuracy that should set them has not
been measured yet.

\section{Behaviour}
\begin{itemize}
   \item \textbf{Binary angle.} $a_i$ is an unsigned fraction of a full circle and wraps at
         $2^{A\_WIDTH}$. A positive angle rotates counter-clockwise.
   \item \textbf{Quadrant step.} The top two bits of $a_i$ select a rotation by 0, 90, 180 or
         270 degrees through swap and negate in the first stage. The remaining
         $A\_WIDTH-2$ bits are the residual angle.
   \item \textbf{Iterations.} Every sample passes all $N\_ITER$ stages, with no early stop.
         Stage $i$ computes $x' = x \mp (y \gg i)$, $y' = y \pm (x \gg i)$ and
         adds or subtracts $\mathrm{atan}(2^{-i})$ from the residual, with the sign taken from the
         residual. $\gg$ is an arithmetic shift that truncates.
         The table entry is
         $\mathrm{round}( \mathrm{atan}(2^{-i}) / (2\pi) \cdot 2^{TABLE\_WIDTH} )$.
   \item \textbf{Word width.} Internally $XY\_WIDTH + GUARD\_BITS + LSB\_BITS$ bits.
   \item \textbf{Guard bits.} In rotation, and in vectoring as stated at the end of this item, the vector length is at most $\sqrt{2} \cdot 2^{XY\_WIDTH-1}$ and the
         CORDIC gain is 1.647, so the intermediate value stays below $2^{XY\_WIDTH+1}$ and two guard
         bits are enough. An exhaustive check over all $x$, $y$ and angle values compared the output with
         two guard bits against the output with eight, for $XY\_WIDTH$ = 5 to 8, $A\_WIDTH$ = 6 and 8,
         $N\_ITER$ = 8, 16, 32, $RM\_GAIN$ = 0, 4, 12 and $LSB\_BITS$ = 0, 1, 4. The word does not wrap
         from $XY\_WIDTH = 6$. At $XY\_WIDTH = 5$ it wraps only at $N\_ITER = 32$ with $LSB\_BITS = 0$,
         and at 2 to 4 bits it wraps in some of the configurations (30, 18 and 12 of 54). Wider $XY\_WIDTH$ and $A\_WIDTH$ are not covered by the
         exhaustive check. Vectoring has the same bound. The same exhaustive check over all $x$ and $y$
         (the angle does not enter), with the parameter sets above, finds no wrapping at $XY\_WIDTH$ = 6 to 8;
         at $XY\_WIDTH = 5$ it wraps in 6 of 54 configurations, all with $N\_ITER = 32$ and $LSB\_BITS = 0$.
         $XY\_WIDTH$ 2 to 4 is not checked for vectoring.
   \item \textbf{Vectoring.} With $VECTORING$ = true the core computes magnitude and angle of
         $(x_i, y_i)$ and ignores $a_i$. $x_o$ is the gain times $\sqrt{x^2+y^2}$, $y_o$ is a
         residual near zero and $a_o$ is $\mathrm{atan2}(y, x)$ as an unsigned binary angle in
         $[0, 2^{A\_WIDTH})$, the format of $a_i$.
         \begin{itemize}
            \item The first stage keeps a vector with $x \ge 0$ and sets the angle to 0. For $x < 0$, $y \ge 0$ it
                  gives $(y, -x)$ with the angle $2^{A\_WIDTH-2}$ (90 degrees); for $x < 0$, $y < 0$ it gives
                  $(-y, x)$ with the angle $3 \cdot 2^{A\_WIDTH-2}$ (270 degrees). The vector is then in the
                  right half-plane, inside the convergence range of $\pm 99.9$ degrees.
            \item Stage $i$ takes its direction from the sign bit of $y$. For $y \ge 0$:
                  $x' = x + (y \gg i)$, $y' = y - (x \gg i)$, $a' = a + \mathrm{atan}_i$; for $y < 0$:
                  $x' = x - (y \gg i)$, $y' = y + (x \gg i)$, $a' = a - \mathrm{atan}_i$. $\mathrm{atan}_i$ is
                  the table entry truncated to $A\_WIDTH$ bits; the angle accumulator has $A\_WIDTH$ bits and wraps
                  modulo $2^{A\_WIDTH}$.
            \item Gain removal acts on $x$ and $y$; the angle is delayed through those stages. Latency,
                  handshake, \signal{en} and reset are those of rotation.
            \item The input $(0, 0)$ gives $x_o = y_o = 0$ and an angle that model and RTL agree on and that
                  carries no meaning.
            \item An unconnected $a_o$ lets synthesis remove the angle registers (table \ref{tab:pipe_synth}).
         \end{itemize}
   \item \textbf{Gain removal.} One pipeline stage per term, with the shift-and-add terms of
         \signal{cordic\_int\_rm\_gain()} of the iterative model. $RM\_GAIN = 0$ leaves the gain 1.647.
         The gain after removal is not exactly 1; it is listed in table \ref{tab:pipe_acc}.
   \item \textbf{Latency.} $1 + N\_ITER + RM\_GAIN$ clocks: the quadrant stage, the iterations
         and the gain stages.
   \item \textbf{Reset.} Synchronous, on the valid chain only. The data registers have no reset.
         The data at $x_o$, $y_o$ and $a_o$ is defined only while \signal{valid\_o} = '1'. The valid chain has
         no initial value, so \signal{valid\_o} is 'U' in simulation until the first rising edge with
         \signal{rst} = '1'. One clock of \signal{rst} with a running clock is enough, whatever \signal{en} is.
\end{itemize}

\subsection{Valid and clock-enable}
\signal{valid\_o} is \signal{valid\_i} delayed by the latency. With \signal{en} = '0' the whole pipeline, including the
valid chain, holds its state, so \signal{valid\_o} and the data stay unchanged. A sample is
therefore taken at an edge where \signal{valid\_o} = '1' \emph{and} \signal{en} = '1': a consumer qualifies
\signal{valid\_o} with \signal{en}. Otherwise it takes the same sample at every frozen edge.
\signal{rst} wins over \signal{en} on the valid chain.

\signal{valid\_i} is taken only at an edge with \signal{en} = '1', so a producer holds its sample
while \signal{en} = '0'. A delay line for sideband data, such as RAM addresses, uses the same
\signal{en}, because the latency counts enabled clocks.

\section{C-Model and Python Wrapper}
\file{cordic\_pipeline.c} is plain C99 with no global state and no \file{mex} dependency.
\signal{cordic\_pipeline()} rotates one sample and \signal{cordic\_pipeline\_array()} an array.
\signal{cordic\_pipeline\_vec( cfg, x\_i, y\_i, x\_o, y\_o, a\_o )} and \signal{cordic\_pipeline\_vec\_array()}
do the same for vectoring; \signal{cordic\_pipeline\_cfg\_t} is the same for both. All four return an error code for an invalid configuration or for $x_i$, $y_i$ outside the signed
$XY\_WIDTH$ range. The model builds with \texttt{-Wall -Wextra -Werror} and, as a
check, with \texttt{-std=c99 -pedantic}.
\file{cordic\_pipeline.py} loads the shared library with \file{ctypes} and has \file{numpy}
arrays at its interface; all arithmetic stays in C. The path of the library is an argument of
\signal{CordicPipeline}. \signal{rotate( x, y, a )} returns $(x_o, y_o)$ and
\signal{vector( x, y )} returns $(x_o, y_o, a_o)$. The angle table and the gain-removal terms are written by
\file{gen\_tables.py} into \file{cordic\_pipeline\_tables.h} and \file{cordic\_pipeline\_pkg.vhd};
\texttt{gen\_tables.py --check} exits with 1 if either file is not current.

\section{Verification}
\begin{itemize}
   \item \texttt{test/model/run.py --out <dir>}: 43 unit tests (12 of them for vectoring) of the model, the wrapper and the
         tables, a sweep of the model against the exact rotation, a sweep against the exact vectoring, and the vectoring
         tests of section \ref{sec:pipe_acc}. It takes about 40 s and ends with \texttt{OK} and
         writes \file{sweep.csv} and \file{sweep\_vec.csv}.
   \item \texttt{test/pipeline/run.py -o <dir>}: VUnit with Questa. A \texttt{pre\_config} hook
         runs the model and writes stimulus and expected data, and the testbench compares
         $x_o$, $y_o$ and $a_o$ exactly for each valid sample; in rotation sets $a_o$ must be zero.
         It runs 7 tests in 16 generic sets (112 tests; 11 rotation sets and 5 vectoring sets, among them
         $12,18,2,0,10,0$, a stand-in for the magnitude CORDIC of gnssr's fast acquisition) in about 116 s with 8 threads: consecutive samples, gaps, \signal{en} toggling, the measured latency, and a
         reset in the stream and while frozen. \texttt{test\_reset\_resume} holds \signal{valid\_i} = '1'
         through a one-clock reset while streaming and requires \signal{valid\_o} = '0' afterwards, so that the
         sample offered at the reset edge never appears. It then reopens the expected file, streams the whole
         stimulus again from its start with random \signal{en}, and requires every sample to compare exactly and
         exactly $N\_SAMPLES$ to arrive. The vectors include all four full-scale corners with
         $-2^{XY\_WIDTH-1}$, zero, and the quadrant boundaries with one LSB on either side.
         The vectoring stimulus carries random $a_i$, which must not change the result, and the directed
         vectoring vectors (corners, $(0, 0)$, axis points, points next to the half-plane boundary).
         The sets include $RM\_GAIN$ = 0 and 12, $LSB\_BITS$ = 0 and above, and two sets in which the word
         wraps. It ends with \texttt{All passed!}.
\end{itemize}

\section{Accuracy}
\label{sec:pipe_acc}
\file{sweep.csv} lists $XY\_WIDTH$ in 12, 14, 16, $A\_WIDTH$ in 16 to 24, $LSB\_BITS$ in 0 to 6,
$N\_ITER$ in 8 to 20 and $RM\_GAIN$ in 0 to 12 (3360 rows). Each row has 20000 uniform
full-scale inputs. The error is the magnitude $|dx + i\,dy|$ in output LSB against the exact
rotation multiplied by the nominal gain of the configuration. The gain error itself is
therefore not in the numbers; the column \emph{gain} is the nominal gain.
\begin{table}[htbp]
   \center\small
   \begin{tabular}{@{}lllll@{}}
      \rowcolor{tableheadcolor}\textbf{XY,A,LSB,N,RM} & \textbf{Latency} & \textbf{Gain} & \textbf{rms} & \textbf{max} \\
      16,18,0,16,0   & 17 & 1.646760 & 3.72 & 18.41 \\\midrule
      16,18,0,16,4 (default) & 21 & 0.997062 & 2.44 & 12.16 \\\midrule
      16,18,0,16,12  & 29 & 1.000000 & 6.29 & 15.76 \\\midrule
      16,18,2,16,4   & 21 & 0.997062 & 1.59 & 7.50 \\\midrule
      14,18,0,10,4   & 15 & 0.997061 & 7.69 & 23.26 \\\midrule
      16,24,6,20,12  & 33 & 1.000000 & 0.74 & 1.39 \\
      \bottomrule
   \end{tabular}
   \caption{Error against the exact rotation, in LSB, from the sweep of \texttt{test/model/run.py}}
   \label{tab:pipe_acc}
\end{table}
$RM\_GAIN = 4$ leaves a gain of 0.997062, so an input of full scale is 0.3 \% short of the
exact length on top of the listed error. In the sweep the error does not fall monotonically with $RM\_GAIN$: at 16,18,0,16 the rms is 1.73 LSB for $RM\_GAIN$ = 2 and 6.29 LSB for 12.
A measurement with 200000 samples at $XY\_WIDTH$ = 16, $A\_WIDTH$ = 18, $LSB\_BITS$ = 0, $N\_ITER$ = 16 gives a mean
error per component of +0.27, +1.28 and +3.27 LSB at $RM\_GAIN$ = 4, 6 and 12, with an rms of 2.44, 3.10 and 6.29 LSB;
with the mean removed the rms is 2.41, 2.52 and 4.25 LSB. At $LSB\_BITS$ = 2 and $RM\_GAIN$ = 12 the offset falls to
0.45 LSB. Confirmed: a growing DC offset accounts for part of the rise. Inferred, not measured: each truncating term
has a bias, and the terms at high $RM\_GAIN$ are mostly negative. Above about $RM\_GAIN$ = 4, more terms therefore need
$LSB\_BITS$.

\subsection{Vectoring}
The model tests compare vectoring with the exact magnitude times the nominal gain and with
$\mathrm{atan2}$. With $G$ the nominal gain, $M = |x + i\,y|$, $e = 2^{-LSB\_BITS}$ and
$r = \mathrm{atan}(2^{-(N\_ITER-1)})$ the limits are derived from the algorithm:
\[
\begin{array}{l}
   |x_o - G M| \le (N\_ITER + RM\_GAIN)\,e + 1 + G M (1 - \cos r),\\[2pt]
   |y_o| \le G M \sin r + (N\_ITER + RM\_GAIN)\,e + 1,\\[2pt]
   |\Delta a| \le \left(r + \frac{(N\_ITER+2)\,e}{g_N M}\right)\frac{2^{A\_WIDTH}}{2\pi} + N\_ITER + 1
   \quad\mbox{for } M \ge 2^{XY\_WIDTH-3},
\end{array}
\]
with $\Delta a$ in angle LSB, wrapped, and $g_N$ the CORDIC gain of $N\_ITER$ stages. The tests cover the axes,
the half-plane boundary, the corners with $-2^{XY\_WIDTH-1}$, $(0, 0)$, $x_o \ge 0$, the fitted gain and
the exhaustive guard-bit check. At the half-plane boundary the angle is checked at $(-s, y)$ with
$s = 2^{XY\_WIDTH-2}$ and $y$ in $0$, $-1$, $\pm s$, which covers both sides of the switch of the start angle between
90 and 270 degrees; the directed pairs at $x = -1$ have a length below $2^{XY\_WIDTH-3}$ and get the magnitude and
$y_o$ checks only.

\file{sweep\_vec.csv} has the same grid as \file{sweep.csv} (3360 rows) and the same 20000 samples per row, with the
columns of the rotation sweep and the rms and maximum of $|x_o - G M|$ in output LSB and of the angle error in angle
LSB, the angle error over the samples with $M \ge 2^{XY\_WIDTH-3}$.

Measured with 20000 uniform full-scale samples (seed 12345) and the wrapper function \signal{vector};
the error of $x_o$ is $x_o - G M$ in output LSB, the angle error is wrapped and in angle LSB
(1 LSB $= 360 / 2^{18} = 0.00137$ degrees at $A\_WIDTH = 18$):
\begin{table}[htbp]
   \center\small
   \begin{tabular}{@{}lll@{}}
      \rowcolor{tableheadcolor}\textbf{XY,A,GUARD,LSB,N,RM} & \textbf{16,18,2,0,16,4} & \textbf{12,18,2,0,10,0 (stand-in)} \\
      nominal gain $G$ & 0.997062 & 1.646759 \\\midrule
      fitted gain & 0.997143 & 1.647499 \\\midrule
      $x_o - GM$ rms / max / mean & 2.70 / 7.14 / +2.35 & 1.94 / 6.34 / +1.38 \\\midrule
      $y_o$ rms / max & 0.34 / 3 & 3.20 / 11 \\\midrule
      angle, $M \ge 2^{XY-3}$, rms / max & 2.05 / 15.2 & 52.3 / 231 \\\midrule
      angle, all samples, rms / max & 3.70 / 216 & 65.5 / 3142 \\
      \bottomrule
   \end{tabular}
   \caption{Vectoring error against the exact magnitude and angle; 12,18,2,0,10,0 is a stand-in for the magnitude CORDIC of gnssr's fast acquisition}
   \label{tab:pipe_vec_acc}
\end{table}
The positive mean error is measured; its cause is not investigated. The angle error at $N\_ITER = 10$ is
dominated by the residual $\mathrm{atan}(2^{-9})$, which is 81.5 angle LSB at $A\_WIDTH = 18$ (inferred from the
bound, not isolated by a measurement). The large maxima over all samples belong to short vectors, where input
quantisation limits the angle.

\section{Synthesis}
Quartus 23.1, Cyclone V 5CSEBA6U23I7, slow 1100 mV 100 C and -40 C models, a scratch project with a 5 ns
clock constraint (3 ns for the vectoring rows) and no pin assignments. The Fmax is therefore indicative only.
\begin{table}[htbp]
   \center\small
   \begin{tabular}{@{}lllllll@{}}
      \rowcolor{tableheadcolor}\textbf{Generics} & \textbf{ALMs} & \textbf{Regs} & \textbf{DSP} & \textbf{RAM} & \textbf{Fmax 100 C} & \textbf{Fmax -40 C} \\
      16,18,2,0,16,4 (default) & 524 & 1068 & 0 & 0 & 252.53 MHz & 243.78 MHz \\\midrule
      14,18,2,0,10,4 (stand-in) & 327 & 676  & 0 & 0 & 276.32 MHz & 267.31 MHz \\\midrule
      12,18,2,0,10,0 vec, a\_o used & 243 & 453 & 0 & 0 & 261.92 MHz & 251.57 MHz \\\midrule
      12,18,2,0,10,0 vec, a\_o open & 177 & 317 & 0 & 0 & 282.97 MHz & 274.95 MHz \\
      \bottomrule
   \end{tabular}
   \caption{Synthesis results}
   \label{tab:pipe_synth}
\end{table}
The generics are XY\_WIDTH, A\_WIDTH, GUARD\_BITS, LSB\_BITS, N\_ITER, RM\_GAIN; RAM is block memory bits. The second set is a stand-in, because the parameters of the FFT that uses the core are open.
The rows marked vec are the stand-in for the magnitude CORDIC of gnssr's fast acquisition, use \texttt{VECTORING} = true and have a top level that exposes only $x_o$, $y_o$, \signal{valid\_o} and, in the first row,
$a_o$; with $a_o$ open Quartus removes 136 registers and 66 ALMs.

\fussy
\chapter{Performance}
The performance of the pipelined core is in chapter \ref{chap:pipeline}. The performance of the iterative core is evaluated through the matlab script \file{cordic\_iterative\_test.m}.


\chapter{Open issues and future work}

   \section{Next steps}
   \begin{itemize}
      \item Add YAC to a FPGA based System-on-Chip to prove 
            FPGA feasibility
      \item Performance optimization
         \item circuit optimizations
         \item numerical optimizations
   \end{itemize}

   \section{Future plans}
   \begin{itemize}
      \item Hyperbolic range extension
      \item Floating point CORDIC 
   \end{itemize}



%%%%%%%%%% BIBLIOGRAPHY %%%%%%%%%%
\begin{thebibliography}{9}

\bibitem{survey}
  Andraka, Ray; \textit{A survey of CORDIC algorithms for FPGA based computers}

\bibitem{vlsi}
   Hu, Yu Hen; \textit{CORDIC-Based VLSI Architectures for Digital Signal Processing}

\bibitem{wikibook}
  CORDIC on wikibook: \url{http://en.wikibooks.org/wiki/Digital_Circuits/CORDIC}

\bibitem{dawid}
   David, Herbert; Meyr, Heinricht; \textit{CORDIC Algorithms and Architectures} 
   \url{http://www.eecs.berkeley.edu/newton/Classes/EE290sp99/lectures/ee290aSp996_1/cordic_chap24.pdf}
\end{thebibliography}
\end{document}


